\chapter{Basi di un programma OpenCL}\label{chap:opencl_basi}
In questo capitolo si spiegano le basi per utilizzare OpenCL seguendo un codice \texttt{vecinit.c} per l'allocazione di un vettore e l'inizializzazione dei suoi elementi al loro indice. Il codice completo è disponibile nella cartella \texttt{codes/vecinit}.

\section{Boilerplate OpenCL}\label{sec:opencl_basi_device_selection}
La parte iniziale di OpenCL è sempre la stessa, e consiste nella selezione del dispositivo, nella creazione del contesto e della coda dei comandi. Il codice seguente mostra come fare tutto questo:

\begin{code}[caption={Boilerplate OpenCL}, label={code:opencl_basi_boilerplate}]
    cl_platform_id p = select_platform();
    cl_device_id d = select_device(p);
    cl_context ctx = create_context(p, d);
    cl_command_queue que = create_queue(ctx, d);

    cl_program prog = create_program("vecinit.ocl", ctx, d);
\end{code}

\subsection*{Descrizione del codice}
Le varie righe possono essere descritte così:
\begin{itemize}
    \item \texttt{select\_platform()} è una funzione che restituisce la piattaforma OpenCL da utilizzare. In un sistema con più piattaforme, è necessario scegliere quella corretta.
    \item \texttt{select\_device(p)} è una funzione che restituisce il dispositivo OpenCL da utilizzare. In un sistema con più dispositivi, è necessario scegliere quello corretto.
    \item \texttt{create\_context(p, d)} è una funzione che crea un contesto OpenCL per la piattaforma e il dispositivo selezionati. Il contesto è un ambiente in cui vengono eseguiti i kernel OpenCL.
    \item \texttt{create\_queue(ctx, d)} è una funzione che crea una coda dei comandi OpenCL per il contesto e il dispositivo selezionati. La coda dei comandi è utilizzata per inviare i comandi al dispositivo.
    \item \texttt{create\_program("vecinit.ocl", ctx, d)} è una funzione che crea un programma OpenCL a partire da un file sorgente. Il programma contiene i kernel OpenCL che verranno eseguiti sul dispositivo. In questo caso, il file sorgente è \texttt{vecinit.ocl}, che contiene il kernel per inizializzare un vettore.
\end{itemize}

\section{Creazione del buffer device}\label{sec:opencl_basi_buffer}

Il buffer è una regione di memoria allocata sul dispositivo che può essere utilizzata per scambiare dati tra il host e il device. Per creare un buffer, è necessario specificare la dimensione del buffer, le opzioni di accesso, il puntatore ai dati (che generalmente è \texttt{NULL} se non si vuole usare un puntatore host) e un puntatore per un eventuale codice di errore. Il codice seguente mostra come creare un buffer:
\begin{code}[caption={Creazione di un buffer OpenCL}, label={code:opencl_basi_buffer}]
    const size_t memsize = sizeof(int)*nels;
    cl_int err;
    cl_mem d_vec = clCreateBuffer(ctx, CL_MEM_WRITE_ONLY, memsize, NULL, &err);
    ocl_check(err, "clCreateBuffer fallito");
\end{code}

\subsection*{Descrizione del codice}
Le varie righe possono essere descritte così:
\begin{itemize}
    \item \texttt{memsize} rappresenta la dimensione del buffer espressa in byte. In questo caso, il buffer contiene \texttt{n} elementi di tipo \texttt{int}, quindi la dimensione totale è pari a \texttt{sizeof(int) * n}.
    
    \item \texttt{clCreateBuffer(ctx, CL\_MEM\_WRITE\_ONLY, memsize, NULL, \&err)} è la funzione che crea il buffer OpenCL. I parametri sono:
    \begin{itemize}
        \item \texttt{ctx} è il contesto OpenCL in cui viene creato il buffer.
        
        \item \texttt{CL\_MEM\_WRITE\_ONLY} è un flag che specifica i permessi di accesso al buffer dal punto di vista del kernel. In particolare, indica che il buffer può essere scritto dal kernel, ma non letto. Altri flag comuni sono \texttt{CL\_MEM\_READ\_ONLY} e \texttt{CL\_MEM\_READ\_WRITE}.
        
        \item \texttt{memsize} indica la dimensione del buffer in byte.
        
        \item \texttt{NULL} indica che non viene fornito alcun puntatore a memoria host iniziale. L'allocazione del buffer è quindi completamente gestita dal runtime OpenCL.
        
        \item \texttt{\&err} è un puntatore a una variabile in cui viene salvato il codice di errore restituito dalla funzione. Se la creazione del buffer ha successo, \texttt{err} assume valore \texttt{CL\_SUCCESS}.
    \end{itemize}
    
    \item \texttt{ocl\_check(err, "clCreateBuffer fallito")} è una funzione di utilità che verifica il valore di \texttt{err}. Se il codice di errore è diverso da \texttt{CL\_SUCCESS}, l’errore viene segnalato e l’esecuzione del programma viene interrotta.
\end{itemize}

\section{Creazione del kernel}\label{sec:opencl_basi_kernel}
Il kernel è la funzione che viene eseguita sul dispositivo. Per creare un kernel, è necessario specificare il nome del kernel e il programma da cui è stato creato. Il codice seguente mostra come creare un kernel:
\begin{code}[caption={Creazione di un kernel OpenCL}, label={code:opencl_basi_kernel}]
	const char *kname = "init_k";
	cl_kernel init_k = clCreateKernel(prog, kname, &err);
	ocl_check(err, "clCreateKernel %s fallito", kname);
\end{code}

\subsection*{Descrizione del codice}
Le varie righe possono essere descritte così:
\begin{itemize}
    \item \texttt{kname} è una stringa che contiene il nome del kernel da creare. In questo caso, il kernel si chiama \texttt{init\_k}.
    \item \texttt{clCreateKernel(prog, kname, \&err)} è la funzione che crea il kernel OpenCL. I parametri sono:
    \begin{itemize}
        \item \texttt{prog} è il programma OpenCL da cui viene creato il kernel. Il programma deve essere stato compilato con successo prima di poter creare un kernel da esso.
        \item \texttt{kname} è il nome del kernel da creare. Questo nome deve corrispondere a una funzione definita nel file sorgente del programma OpenCL.
        \item \texttt{\&err} è un puntatore a una variabile in cui viene salvato il codice di errore restituito dalla funzione. Se la creazione del kernel ha successo, \texttt{err} assume valore \texttt{CL\_SUCCESS}.
    \end{itemize}
    \item \texttt{ocl\_check(err, "clCreateKernel \%s fallito", kname)} è la solita funzione di errore che verifica il valore di \texttt{err} e segnala eventuali errori nella creazione del kernel.
\end{itemize}

\section{Esecuzione del kernel}\label{sec:opencl_basi_execution}
Nella funzione \texttt{init} del sorgente C viene assegnato ad ogni elemento del vettore il suo indice. In questa funzione viene eseguito il kernel \texttt{init\_k} con un certo numero di work-item. Il codice seguente mostra come eseguire il kernel:
\begin{code}[caption={Esecuzione di un kernel OpenCL}, label={code:opencl_basi_execution}]
  cl_int err;
  cl_event init_evt;

  const size_t gws[1] = { n };

  cl_uint arg_index = 0;
  err = clSetKernelArg(init_k, arg_index, sizeof(cl_mem), &vec);
  ocl_check(err, "init_k set kernel arg %u", arg_index);

  err = clEnqueueNDRangeKernel(que, init_k,
    1,    /* dimensionalita' griglia: nel nostro caso 1D */
    NULL, /* global work offset */
    gws,  /* global work size: quanti elementi totali lanciare */
    NULL, /* local work size: come raggruppare gli elementi */
    0, NULL, /* wait list: non aspettiamo nulla */
    &init_evt /* evento associato al lancio del kernel */);

  ocl_check(err, "enqueue kernel init_k");
\end{code}

\subsection*{Descrizione del codice}
Le varie righe possono essere descritte così:
\begin{itemize}
    \item \texttt{cl\_event init\_evt} è una variabile di tipo \texttt{cl\_event} utilizzata per tenere traccia dell'esecuzione di un comando sulla command queue. In questo caso, \texttt{init\_evt} sarà associato al lancio del kernel \texttt{init\_k} e potrà essere usato successivamente per sincronizzazione o misure temporali.
    
    \item \texttt{gws} (global work size) è un array che contiene la dimensione globale del dominio di esecuzione del kernel. Nel caso monodimensionale, \texttt{gws[0]} rappresenta il numero totale di work-item. In questo caso vale \texttt{n}.
    
    \item \texttt{clSetKernelArg(init\_k, arg\_index, sizeof(cl\_mem), \&vec)} è la funzione che imposta un argomento del kernel. I parametri sono:
    \begin{itemize}
        \item \texttt{init\_k} è il kernel su cui si vuole impostare l'argomento.
        
        \item \texttt{arg\_index} è l'indice dell'argomento. Gli argomenti sono indicizzati a partire da 0.
        
        \item \texttt{sizeof(cl\_mem)} è la dimensione dell'argomento in byte.
        
        \item \texttt{\&vec} è un puntatore al valore dell'argomento. In questo caso, \texttt{vec} è un buffer OpenCL passato al kernel.
    \end{itemize}
    
    \item \texttt{clEnqueueNDRangeKernel(que, init\_k, 1, NULL, gws, NULL, 0, NULL, \&init\_evt)} è la funzione che accoda il kernel per l'esecuzione sulla command queue. I parametri sono:
    \begin{itemize}
        \item \texttt{que} è la command queue associata al dispositivo su cui verrà eseguito il kernel.
        
        \item \texttt{init\_k} è il kernel da eseguire.
        
        \item \texttt{1} indica che il dominio di esecuzione è monodimensionale.
        
        \item \texttt{NULL} è il global work offset. Non essendo specificato, l'esecuzione parte dall'indice globale 0.
        
        \item \texttt{gws} è il global work size, cioè il numero totale di work-item da eseguire.
        
        \item \texttt{NULL} è il local work size. Non essendo specificato, la sua scelta è demandata al runtime OpenCL.
        
        \item \texttt{0} indica che non ci sono eventi da attendere prima dell'esecuzione.
        
        \item \texttt{NULL} è la lista degli eventi da attendere, che in questo caso è vuota.
        
        \item \texttt{\&init\_evt} è un puntatore alla variabile evento in cui viene salvato l'evento associato al lancio del kernel.
    \end{itemize}
\end{itemize}

\section{Recupero dei dati}\label{sec:opencl_basi_retrieval}
Dopo l'esecuzione del kernel, è necessario recuperare i dati dal dispositivo per renderli disponibili sull'host. Per fare questo si utilizza la funzione \texttt{clEnqueueReadBuffer}, che legge i dati da un buffer OpenCL e li copia in una regione di memoria host. Il codice seguente mostra come fare:
\begin{code}[caption={Recupero dei dati da un buffer OpenCL}, label={code:opencl_basi_readbuffer}]
    int *h_vec = malloc(sizeof(int) * n);
    if (!h_vec) {
        fprintf(stderr, "allocazione fallita\n");
        exit(3);
    }

    cl_event wait_list[1] = { init_evt };
    cl_event read_evt;

    err = clEnqueueReadBuffer(
        que,           /* coda dei comandi */
        d_vec,         /* buffer da cui leggere */
        CL_TRUE,       /* chiamata bloccante per l'host */
        0, memsize,    /* offset e numero di byte da copiare */
        h_vec,         /* destinazione in memoria host */
        1, wait_list,  /* evento da attendere prima della lettura */
        &read_evt
    );    /* evento associato alla lettura */

    ocl_check(err, "read buffer d_vec");
\end{code}

\subsection*{Descrizione del codice}
Le varie righe possono essere descritte così:
\begin{itemize}
    \item \texttt{int *h\_vec = malloc(sizeof(int) * n)} alloca memoria sul lato host per contenere i dati letti dal buffer OpenCL. In particolare, vengono allocati \texttt{n} elementi di tipo \texttt{int}.
    
    \item \texttt{if (!h\_vec)} è la verifica che l’allocazione sia andata a buon fine. In caso contrario, viene stampato un messaggio di errore e l’esecuzione del programma viene terminata.
    
    \item \texttt{cl\_event wait\_list[1] = \{ init\_evt \}} è un array di eventi che rappresenta la lista di comandi che devono essere completati prima di poter eseguire la lettura. In questo caso, si aspetta il completamento del kernel associato all’evento \texttt{init\_evt}.
    
    \item \texttt{cl\_event read\_evt} è una variabile che conterrà l’evento associato all’operazione di lettura dal buffer.
    
    \item \texttt{clEnqueueReadBuffer(...)} è la funzione che accoda il comando di lettura del buffer dalla memoria del device alla memoria host. I parametri sono:
    \begin{itemize}
        \item \texttt{que} è la command queue su cui viene accodata l’operazione.
        
        \item \texttt{d\_vec} è il buffer OpenCL da cui vengono letti i dati.
        
        \item \texttt{CL\_TRUE} indica che la chiamata è bloccante per l’host, cioè la funzione ritorna solo quando i dati sono stati completamente copiati.
        
        \item \texttt{0} è l’offset in byte da cui iniziare la lettura.
        
        \item \texttt{memsize} è il numero di byte da copiare.
        
        \item \texttt{h\_vec} è il puntatore alla memoria host in cui vengono copiati i dati.
        
        \item \texttt{1} è il numero di eventi nella wait list.
        
        \item \texttt{wait\_list} è la lista di eventi che devono essere completati prima di eseguire la lettura.
        
        \item \texttt{\&read\_evt} è un puntatore alla variabile in cui viene salvato l’evento associato all’operazione di lettura.
    \end{itemize}
    
    \item \texttt{ocl\_check(err, "read buffer d\_vec")} verifica che l’operazione sia andata a buon fine. In caso di errore, esso viene segnalato e l’esecuzione viene interrotta.
\end{itemize}